% Written answer to every assignment question, in sentences. Used by the section reports
% (\answer{\Ans...} under each question box) and by the appendix slides (one per question).
% Numbers are macros from data/values.tex (Reports/collect_values.py).
\newcommand{\AnsOneOne}{%
Our cable is the NKT 33 kV single-core XLPE-AL cable \(1\times185/50\) mm\textsuperscript2: a solid
aluminium conductor (radius 7.4 mm), a 0.4 mm semiconducting conductor screen, 8 mm of XLPE, a 0.4 mm
semiconducting insulation screen, a 50 mm\textsuperscript2 copper wire screen with a copper tape, and a PE sheath with
an outer diameter of 43 mm. It is rated 33 kV between phases (19 kV to ground) and 350 A continuously at a conductor
temperature of 90\,\(^\circ\)C, for three cables in touching trefoil laid directly in ground 1.0 m deep, in soil
of 1.2 K\,m/W at 15\,\(^\circ\)C, with the screens bonded at both ends. The tables give the radius of every layer
with its basis, and the electromagnetic and thermal properties of every material with its source.}
\newcommand{\AnsOneTwo}{%
We take the environment of the datasheet rating: homogeneous soil with a thermal resistivity of
1.2 K\,m/W at 15\,\(^\circ\)C, the trefoil centre 1.0 m below the ground surface, and air at 15\,\(^\circ\)C above
it. The electrical soil properties (\(\varepsilon_r=10\), \(\sigma=0.01\) S/m, from the course) hardly matter:
the grounded screens shield the electric field from the soil, and the soil skin depth of about 700 m makes its eddy
currents negligible. Like the datasheet rating, we neglect soil drying and layering.}
\newcommand{\AnsOneThree}{%
We assume both-end (solid) bonding with ideal earth connections, because that is the bonding of the
NKT 350 A rating: circulating screen currents are allowed. Single-point bonding, with no net screen current, is
studied as a loss-reducing alternative in Q3.11 and Section~5. The cable has no armour and no pipe.}
\newcommand{\AnsOneFour}{%
We impose the phase-to-ground voltage \(\hat V=\sqrt2\cdot33/\sqrt3=26.944\) kV and the phase current
\(\hat I=\sqrt2\cdot350=494.97\) A as peak phasors at 50 Hz, in positive sequence: phase A at \(0^\circ\), B at
\(-120^\circ\) and C at \(+120^\circ\). The voltages drive the electric model of Section~2 and the currents the
magnetic model of Section~3.}
\newcommand{\AnsTwoOne}{%
We solve \(-\nabla\cdot\big((\sigma+j\omega\varepsilon)\nabla\hat v\big)=0\) for the complex potential in
the three XLPE annuli only, with linear triangles. The boundary conditions are \(\hat v=\hat V_k\) on each conductor
screen and \(\hat v=0\) on each core screen (Dirichlet conditions of the course function space). We truncate the domain at the core screens because a continuous grounded screen fixes the potential outside
it to zero. Only the XLPE (its radii and permittivity) and the two screen potentials influence the insulation
field; the conductor, the copper screen, the fillers, the sheath, the soil and the other two cables do not.}
\newcommand{\AnsTwoTwo}{%
Yes. Leaving out everything outside the grounded screens is exact. Treating the semiconducting
layers as ideal electrodes costs a voltage drop of about \AnaDrop\ mV, a fraction \(3\cdot10^{-7}\) of the applied
voltage. Replacing touching jackets by a 0.5 mm gap changes the maximum field by \dEtouch\,\% and the capacitance
by less than \dCtouch\,\%. The one detail that must be kept is the semiconductor thickness: putting the voltage on
the aluminium surface instead of on the semicon would change the capacitance by \dCbare\,\%.}
\newcommand{\AnsTwoThree}{%
The capacitance is converged already on the coarsest mesh (changes of order \(10^{-5}\)). The maximum
field converges from below: it is \dEgap\,\% from the analytical value on the default mesh and \dEfine\,\% on the
fine one (\Dofsgapfine\ complex unknowns). With linear elements the field is constant per triangle, so the error is
concentrated at the conductor screen, where \(E\propto1/r\) is steepest: local refinement is needed there (and
around a defect), not in the bulk of the insulation.}
\newcommand{\AnsTwoFour}{%
The datasheet gives no admissible field, so we use the design limits of the Eskom XLPE cable
specification: 8 kV/mm RMS at the conductor screen and 4 kV/mm RMS at the core screen. The healthy cable has
\AnaEirms\ kV/mm and \AnaEorms\ kV/mm RMS at these two screens, 43\,\% of each limit, so it is well within them.}
\newcommand{\AnsTwoFive}{%
We add an air-filled void (a gas cavity) of radius 0.25 mm with \(\varepsilon_r=1\), centred at
\(r=8.8\) mm so that its edge is 0.75 mm from the conductor screen, where the stress is highest. In a 2D cross-section the void is an infinitely long channel
along the cable; its field enhancement (\(\AnaVoidFactor\,E_0\)) is larger than that of a real spherical void (\(\AnaSphereFactor\,E_0\)), so
the 2D result is conservative. Averaged over the void, the field rises to \Vmeandefectfine\ kV/mm
(\(\VoidRatio E_0\)), while the field in the surrounding XLPE stays close to the healthy one. The cable does not
withstand this defect in the long run: the void field exceeds the breakdown strength of air (about \AnaEb\ kV/mm),
so partial discharges start at rated voltage and slowly erode the insulation until it breaks down.}
\newcommand{\AnsTwoSix}{%
At the conductor screen the displacement current density \(\omega\varepsilon|\hat E|\) is
\AnaJd\ A/m\textsuperscript2 and the conduction current density \(\sigma|\hat E|\) is \AnaJc\ A/m\textsuperscript2;
their ratio \(\sigma/(\omega\varepsilon)=\AnaRatio\) is the same everywhere in the XLPE. Current conservation is
satisfied: the FE equations are the weak form of \(\nabla\cdot(\hat{\mathbf J}+j\omega\hat{\mathbf D})=0\), one current
balance per node, so the current entering at the conductor screen, \IchFE\ A/km per phase (\(\omega CV\)), leaves
at the core screen. Displacement current dominates in the insulators (XLPE, PE and air), where
\(\sigma\ll\omega\varepsilon\); conduction current dominates in the semicons, the metals and the soil, where
\(\sigma\gg\omega\varepsilon\).}
\newcommand{\AnsTwoSeven}{%
The FE capacitance is \Cgap\ \(\mu\)F/km per phase. It equals the analytical value
\(2\pi\varepsilon/\ln(r_o/r_i)=\AnaC\) \(\mu\)F/km to \dCana\,\% and lies below the NKT maximum of 0.22
\(\mu\)F/km. The analytical formula assumes concentric, homogeneous and linear insulation and treats the
semiconducting layers as perfect electrodes, so \(r_i\) and \(r_o\) are the radii of the XLPE interfaces (7.8 and
15.8 mm), not of the aluminium or the copper. The capacitance does not vary with frequency: from 5 Hz to 5 kHz it
stays at \Cgap\ \(\mu\)F/km, because in a homogeneous annulus the factor \(\sigma+j\omega\varepsilon\) cancels from
the equation.}
\newcommand{\AnsTwoEight}{%
We propose a thicker insulation: 9 mm instead of 8 mm of XLPE at the same outer diameter. This
changes the peak stress at the conductor screen by \dEthick\,\% (from \Egap\ to \Ethick\ kV/mm peak) and the
capacitance and charging current by \dCthick\,\%, which gives more margin against defects such as the void of Q2.5.}
\newcommand{\AnsThreeOne}{%
We solve \(\nabla\times(\nu\nabla\times\mathbf a)=\mathbf j\) for \(\mathbf a=a_z\mathbf e_z\) in a
\(0.6\times0.6\) m box that contains the three cables and the soil between and around them, with \(a_z=0\) on the
box edge. All cables must be in the domain because the screens do not shield the magnetic field. We truncate at
0.3 m because the balanced phase currents make the field decay like \(1/r^2\): at the box edge it is below 1\,\%
of its peak. The aluminium and the copper are massive regions with \(\mathbf j=-\sigma(j\omega\mathbf a+\mathbf
u_r)\); each phase current is imposed as the net current of its aluminium region (\texttt{Current\_2D}), so that
skin and proximity effects are free.}
\newcommand{\AnsThreeTwo}{%
The skin depth in aluminium at 50 Hz is \AnaDeltaAl\ mm, larger than the 7.4 mm conductor radius,
so skin effect is weak. The massive model is the realistic one, because NKT specifies a solid aluminium conductor in
which eddy currents flow freely. At 50 Hz the two models differ little (conductor resistance \dRcstr\,\%,
inductance \dLstr\,\%), but only the massive model shows the current crowding towards the other phases and the
skin and proximity increase of the resistance (\SkinFE\,\%).}
\newcommand{\AnsThreeThree}{%
Each copper screen is a massive region with its own voltage per unit length \(U\) and net current
\(I\). Both-end bonding imposes \(U=0\) and the circulating current follows from the solution; single-point bonding
imposes \(I=0\) and a standing voltage follows. Using only \(\mathbf j=-j\omega\sigma\mathbf a\) in the screens
silently sets \(U=0\) relative to the arbitrary reference \(a_z=0\) on the truncation boundary, so the net screen
current would depend on the size and shape of the box instead of on a real earthing, and \(I=0\) could not be
imposed at all.}
\newcommand{\AnsThreeFour}{%
Yes, within limits. Replacing touching jackets by a 0.5 mm gap changes the inductance by
\dLtouch\,\% and \(R_{AC}\) by \dRactouch\,\%, and modelling the copper wires and tape as one ring of the same
area keeps the DC resistance and the screen loop. The main simplification of Section~2 is no longer valid: the
screens do not confine the magnetic field, so all metals, all three cables, their spacing and the soil must be
kept. Only the simplifications of the semicons and fillers remain valid.}
\newcommand{\AnsThreeFive}{%
The mesh targets 0.25 mm around the cables; the element size in the aluminium and copper
is at most \MeshHmetalGap\ mm, below 5\,\% of the skin depth (\AnaDeltaCu\ mm in copper, \AnaDeltaAl\ mm in aluminium), so
the mesh is limited by the geometry, not by the skin depth. Refining from MeshFactor 1 to 0.5 changes
\(R_{AC}\) by \dRacfine\,\%, the inductance by \dLfine\,\% and the screen losses by \dPsfine\,\%. The screen
losses converge slowest because the current varies across the thin copper ring, so that is where to refine.}
\newcommand{\AnsThreeSix}{%
The flux density amplitude is highest, \Bmax\ mT (peak), on the conductor surfaces that face the
centre of the trefoil.
It grows linearly inside each conductor, decreases like \(1/r\) across the insulation, and the field of the
balanced trefoil decays quickly in the soil, to below 1\,\% of its peak at the box edge.}
\newcommand{\AnsThreeSeven}{%
In the aluminium the mean current density amplitude is \Javg\ A/mm\textsuperscript2 (peak) and the maximum is
\Jmax\ A/mm\textsuperscript2, at the surface facing the other phases (proximity effect). With both ends bonded the
screens carry a circulating current of about \Jscreen\ A/mm\textsuperscript2 on average; with single-point bonding
only local eddy currents remain, with zero net current.}
\newcommand{\AnsThreeEight}{%
With both ends bonded, the aluminium conductors dissipate \Pcbase\ kW/km and the copper screens
\Psbase\ kW/km, \Ptbase\ kW/km in total for the three phases. With single-point bonding the screen losses drop to
\Pssingle\ kW/km. There is no armour or pipe, and the soil has no losses because its conductivity is neglected.}
\newcommand{\AnsThreeNine}{%
The AC resistance defined from the total losses, \(R_{AC}=P/(3I^2)\), is \Racbase\ \(\Omega\)/km
at 20\,\(^\circ\)C, against a DC resistance of \Racfdc\ \(\Omega\)/km (NKT: at most 0.164 \(\Omega\)/km). Skin and
proximity effect add \SkinFE\,\% (IEC 60287: \IecYsp\,\%) and the screen losses add \ScreenFE\,\%; IEC 60287 gives
\IecRac\ \(\Omega\)/km in total. The datasheet does not list an AC resistance.}
\newcommand{\AnsThreeTen}{%
For a balanced, symmetric three-phase system we define the positive-sequence inductance per phase
from the stored magnetic energy, \(L=\int\nu|\hat{\mathbf b}|^2dA/(3\hat I^2)\); it includes the internal
inductance and the reaction of the screen currents. It is \Lbase\ mH/km at 50 Hz and it does vary with frequency
(\Lfdc\ mH/km at DC, \Lfhund\ mH/km at 500 Hz), because the induced screen and eddy currents oppose the flux more
at higher frequency. For touching cables the FE gives \Ltouching\ mH/km, against \NktL\ mH/km from the NKT
reactance of 0.12 \(\Omega\)/km (\dLnkt\,\%).}
\newcommand{\AnsThreeEleven}{%
We propose single-point (or cross-) bonding of the screens. It removes the circulating screen
current, so the total losses fall from \Ptbase\ to \Ptsingle\ kW/km (\dPtSingle\,\%) and \(R_{AC}\) from
\Racbase\ to \Racsingle\ \(\Omega\)/km. The price is a standing voltage of \Usingle\ V/km on the open screen end,
which limits the section length and requires sheath voltage limiters.}
\newcommand{\AnsFourOne}{%
We solve steady-state heat conduction, \(-\nabla\cdot(\kappa\nabla T)=p\), with the Joule loss
density \(p=|\hat{\mathbf j}|^2/(2\sigma)=\tfrac12\sigma(T)|j\omega\hat{\mathbf a}+\hat{\mathbf u}_r|^2\) as heat
source in the aluminium and copper. The factor \(\tfrac12\) turns the complex peak phasors into a real,
time-averaged loss density; in GetDP the expression is evaluated on the complex magnetic solution with
\texttt{<a>[...]} and \texttt{SquNorm}, which returns a real number. In the coupled model the conductivity depends
on temperature, \(\sigma=\sigma_{20}/(1+\alpha(T-20))\).}
\newcommand{\AnsFourTwo}{%
The cables are cooled by conduction through the soil to the ground surface, where the heat leaves
by convection and radiation; we model this with \(-\kappa\,\partial_nT=h(T-T_{\rm amb})\) and
\(h=10\) W/(m\(^2\)K) (natural convection about 2, linearised radiation \Hrad, plus some wind). We truncate the
domain 20 m to each side, with zero heat flux, and 20 m deep, with \(T=15\,^\circ\)C. The thermal domain must be
much larger than the electromagnetic one because the three cables form a net heat source whose temperature field
decays only like \(\ln r\), while the field of the balanced currents decays like \(1/r^2\). We use one mesh: the
0.6 m electromagnetic box is a sub-domain of the thermal domain, so no interpolation is needed.}
\newcommand{\AnsFourThree}{%
The ambient temperature is 15\,\(^\circ\)C for the air and the deep soil, as in the NKT rating.
Refining the mesh from MeshFactor 1 to 0.5 changes the hot spot by \dTcFine\ K, and enlarging the domain from 20 to
40 m changes it by \dTcDom\ K, so the default mesh and domain are adequate. A 10 m domain changes the hot spot by
\dTcDomten\ K, because its 15\,\(^\circ\)C bottom then acts as an extra heat sink.}
\newcommand{\AnsFourFour}{%
The hot spot is \Tcnl\,\(^\circ\)C, in the aluminium of the two lower cables (B and C); the top cable
is \dTcAlow\ K cooler because it is closer to the ground surface. The aluminium is nearly isothermal, and most of
the temperature drop is across the XLPE and the soil. At the rated 350 A the hot spot is \dTcNkt\ K from the
90\,\(^\circ\)C limit for XLPE given by NKT: the cable is at its thermal limit, within the uncertainty of our
assumptions.}
\newcommand{\AnsFourFive}{%
The ground above the trefoil is warmed by \dTsurf\ K, to a maximum of \Tsnl\,\(^\circ\)C at the
ground/air interface, and this warming spreads over a few metres. With a Dirichlet condition \(T=15\,^\circ\)C at
the surface, the hot spot drops by \dTcIso\ K to \Tcnliso\,\(^\circ\)C: the isothermal surface is slightly
non-conservative, because it removes the small thermal resistance \(1/h\) of the surface.}
\newcommand{\AnsFourSix}{%
The IEC 60287 calculation gives \IecTheta\,\(^\circ\)C at 350 A and an ampacity of \IecAmp\ A, which
confirms the NKT rating with our material data. Our FE model gives \Tcnl\,\(^\circ\)C (\Tcnliso\,\(^\circ\)C with an
isothermal surface) and an ampacity of \Amptrefoil\ A, so it predicts the admissible conductor temperature within a
few kelvin. The differences come from the convective ground surface (IEC assumes an isothermal one), the empirical
IEC formula for the soil resistance of a trefoil against our full 2D conduction with mutual heating, the eddy losses
included in our screens, and the 0.5 mm gap.}
\newcommand{\AnsFourSeven}{%
Our criterion is the relative increment of the vector of all nodal temperatures in the Euclidean
norm, \(\|T_k-T_{k-1}\|_2/\|T_k\|_2<10^{-6}\). The staggered iteration needs \Nitnl\ iterations, the same number on
the coarse, default and fine meshes, so the convergence does not depend on the mesh density. The convergence
history is a straight line on the logarithmic scale.}
\newcommand{\AnsFourEight}{%
Compared with the linear model (\(\sigma\) at 20\,\(^\circ\)C), the temperature-dependent
conductivity raises the conductor losses by \dPcNL\,\%, the total losses and \(R_{AC}\) by \dPtNL\,\%, and the hot
spot by \dTcNL\ K, because aluminium near 90\,\(^\circ\)C is about \RhoHot\,\% more resistive while the current is
imposed. The screen losses decrease (\dPsNL\,\%), because a more resistive screen carries less circulating
current. The flux density, the inductance (\dLNL\,\%) and the current density barely change, because the currents
are imposed and the materials are non-magnetic.}
\newcommand{\AnsFourNine}{%
No, the linear model is not conservative: it underestimates the losses and the hot spot (by
\dTcNL\ K). Used for rating, it would allow about \AmpLin\ A instead of \Amptrefoil\ A, so the conductor would run
above 90\,\(^\circ\)C and the insulation would age faster. A linear model is only acceptable with \(\sigma\)
evaluated at the maximum operating temperature, which is what IEC 60287 does.}
\newcommand{\AnsFive}{%
For the specified 350 A the best arrangement is the datasheet's touching trefoil: it is the most
compact, it has the lowest inductance and external magnetic field, and it reaches \Tctrefoil\,\(^\circ\)C with both
ends bonded (at its limit) or \Tctrefoilsp\,\(^\circ\)C with single-point bonding. If more current is needed, a
flat formation with 70 mm free distance and single-point (or cross-) bonded screens gives the largest ampacity,
\Ampflatsevsp\ A against \Amptrefoil\ A for the trefoil, at the price of a trench about three times wider and a
higher inductance. With both ends bonded, spacing helps much less (\Ampflatsev\ A), because the larger screen loops
carry more circulating current.}
